% N22 FSE 2027 Industry Track — truthful experimental RESULTS ONLY.
% Exploratory retrospective results; not a finished, accepted, or submission-ready paper.
% Requires: \usepackage{booktabs} in the parent ACM article preamble.
\section{Exploratory Evaluation}
\label{sec:n22-evaluation}

\subsection{Research Questions and Method}
We ask whether completed migration paths that converge to identical relational
schemas and migration ledgers can nevertheless violate application-relevant
data contracts (RQ1), whether contract-aware checks discriminate real defects
from legitimate user modifications better than structural and untyped data
comparisons (RQ2), and whether they add detection power over a \emph{fully
asserted} migration-test baseline at equal replay budget (RQ3).

We used the unmodified migration engines of Django 6.0.0, Django 6.0.1, Django
5.1.8, and OpenCode v2.0.19 (Git revision
\texttt{1fd016ef32286de9489b7b24f1029f52c49a27b3}). Each Django input
was exercised in an independently created SQLite database through the genuine
\texttt{manage.py migrate} command; OpenCode used the actual
\texttt{V1Migration.run} entry point on a new in-memory database following its
real schema migrations. We captured raw business data, schema DDL, migration
ledger, and an application-level or source-projection oracle. Each of 46
distinct case inputs was repeated three times: 138 observations generated in
126 separate operating-system processes. We calculated predictions at the
\emph{case} level, not the repeat level.

Our primary sample consists of two retrospectively selected, independently
reported positive defects from different upstream projects and 13 distinct
negative scenarios. We separately report 30 correlated, legitimate site
configuration variations; one mixed-writer scenario was excluded because the
upstream support policy for continuing V1 writes after a completed V2
checkpoint was not established. The experimental design and labels were
\emph{not} prospectively preregistered or blind; no population prevalence or
unbiased detector sensitivity can be inferred from the sample.

\subsection{Convergent-Structure Failures (RQ1)}
In Django 6.0.0, the \texttt{RenameModel(Song, Son)} migration accidentally
rewrites the permission \texttt{view\_songprofile} held by an existing user
to \texttt{view\_son}. The native \texttt{User.has\_perm} call therefore
returns false after a staged upgrade, whereas granting the same logical
permission after a fresh install succeeds. The staged and fresh databases have
the same final schema checksum and migration-ledger checksum. The regression
was independently reported as Django ticket \#36843 and is absent in the
Django 6.0.1 fix \cite{django36843}.

The second case is the OpenCode V1-to-V2 message conversion at v2.0.19. A
historically valid V1 assistant message without a newly required
\texttt{agent} field is omitted: one of two source messages is absent from the
V2 message table although the migration checkpoint reads
\texttt{\{"phase":"completed"\}}. A control with the current message
shape copies both rows; the schema and migration ledgers match. Independent
issue \#52226 describes the same historical-format incompatibility
\cite{opencode52226}. We do not equate the reporter's production-scale count
with our two-row synthetic reproduction.

\subsection{Differential Comparators (RQ2--RQ3)}
\begin{table}[t]
  \centering
  \caption{Case-level retrospective confusion matrices. Thirteen negative
  inputs exclude the 30 correlated user-edit stress variants. A strong
  baseline already given the same business contracts ties the approach.}
  \label{tab:n22-confusion}
  \begin{tabular}{lrrrr}
    \toprule
    Comparator & TP & FN & FP & TN \\
    \midrule
    Business data contract & 2 & 0 & 0 & 13 \\
    Full baseline + same assertions & 2 & 0 & 0 & 13 \\
    Schema equality only & 0 & 2 & 0 & 13 \\
    Migration ledger only & 0 & 2 & 0 & 13 \\
    Completed status only & 0 & 2 & 0 & 13 \\
    Untyped row difference & 2 & 0 & 5 & 8 \\
    Row-count difference only & 2 & 0 & 4 & 9 \\
    \bottomrule
  \end{tabular}
\end{table}

Both confirmed defects were found by the contract check; no core-negative
input was flagged. Neither checking schema equality, checking migration ledger
equality, nor accepting a successful migration status detected either defect.
Uninterpreted row differences flagged five legitimate cases, including
preserved site edits and path-specific permission rows; row counts alone
flagged four. Across 30 additional correlated user-controlled
\texttt{Site.domain/name} configurations, the contract check flagged none,
whereas the untyped row comparator flagged all 30. These are robustness
observations, \emph{not} additional independent project-level negatives.

The strong baseline receives exactly the same application contracts and
replays exactly the same paths. It discovers both defects with zero false
alarms, so the observed \emph{incremental recall is zero percentage points}.
A more limited baseline using \texttt{django-test-migrations} 1.7.0 did not
expose the Django permission bug because that helper mutes migration signals,
not because developers could not write equivalent assertions. Likewise the
unmodified OpenCode migration-test suite passed 26/26 tests across each of
three runs but did not exercise the missing-\texttt{agent} historical shape.
These results characterize specific test-coverage gaps, rather than proving
superiority over complete existing testing \cite{djangoTestMigrations}.

\subsection{Statistical Uncertainty, Threats, and Decision}
The two defects are two selected historical incidents, not a probability
sample. Even under an optimistic independent-binomial assumption, observing
2/2 detected positives yields a one-sided exact 95\% recall lower bound of
only 22.36\%; observing 0/13 false positives yields an upper bound of 20.58\%.
Project clustering and retrospective case selection invalidate treating these
bounds as population guarantees. We did not evaluate real production datasets,
PostgreSQL/MySQL migrations, organization-specific rollback policy, developer
effort, incident cost, or industrial ROI. Research prior art already includes
fresh-versus-upgrade migration checks \cite{artemisMigrations}.

The original research gate required at least two independent real defects,
at least 15 percentage points of recall gain over the \emph{complete} baseline
at equal replay budget, and at most 5\% false positives. The first condition
is represented in this sample, but the second fails (0 pp gain), and the
small sample cannot support the third with useful precision. Thus, we
\emph{do not claim the proposed industrial-method contribution has been
validated}. The artifact provides reproducible failure characterization and
negative controls; an industry-impact claim requires a prospective,
multi-project deployment evaluation and a measured operational advantage.
